As we have a differential equation that should be integrated, we expect that
the student can solve it in two ways: using an approximation (Way 1) or
calculus (Way 2). \hfill(2 pt finding final scale factor)

\medskip
\noindent\textit{Way 1.}
Using the definition of Hubble constant:
$\dfrac{da}{dt} = Ha$, where $a$ is the scale factor and using the
approximation $\dfrac{da}{dt} \approx \dfrac{\Delta a}{\Delta t}$, and
supposing that the variation of the scale factor and the Hubble constant is
small compared to their values, and using that $aT = \mathrm{const}$:
\begin{equation*}
\frac{\Delta a}{\Delta t} = H_i a_i
\;\Rightarrow\;
\frac{\Delta T}{\Delta t} = -H_i T_i
\;\Rightarrow\;
\Delta t = -\frac{1}{H_i}\,\frac{\Delta T}{T_i} = 5.0\cdot10^{8}~years
\tag*{(7 pt + 1 pt)}
\end{equation*}

\medskip
\noindent\textit{Way 2.}
As we have a matter dominated Universe ($\rho \propto a^{-3}$), without dark
energy, and a density parameter $\Omega_0 = 1$, we have:
\begin{equation*}
\Omega_0 = \frac{8\pi G\rho}{3H^2} = \frac{8\pi G}{3H^2}\,\frac{\zeta}{a^3} = 1
\tag*{(4 pt)}
\end{equation*}
From the Hubble law, $\dfrac{da}{dt} = Ha$:
\[
\left(\frac{da}{dt}\right)^{2} = \frac{8\pi G}{3}\,\frac{\zeta}{a}
\;\Rightarrow\;
\frac{da}{dt} = \sqrt{\frac{8\pi G}{3}\,\frac{\zeta}{a}}
\]
Integrating this equation: \hfill(3 pt)
\[
\frac{2}{3}\left[a_f^{3/2} - a_i^{3/2}\right]
= \sqrt{\frac{8\pi G}{3}\,\zeta}\,\Delta t
\;\Rightarrow\;
\Delta t
= \frac{2}{3}\left[\left(\frac{a_f}{a_i}\right)^{3/2} - 1\right]
  \frac{1}{\sqrt{\dfrac{8\pi G}{3}\,\dfrac{\zeta}{a_i^{3}}}}
= \frac{2}{3}\,\frac{1}{H_i}
  \left[\left(\frac{a_f}{a_i}\right)^{3/2} - 1\right]
\]
Using the fact that $aT = \mathrm{const}$:
\[
a_i T_i = a_f T_f
\;\Rightarrow\;
\frac{a_f}{a_i} = \frac{T_i}{T_f} = \frac{2.73}{2.63} = 1.038
\]
So the time is:
\begin{equation*}
\Delta t = 5.2\cdot10^{8}~years \tag*{(1 pt)}
\end{equation*}
Obs: (7 pt for using correctly that $a \propto t^{2/3}$)
